\chapter*{Annexe A}
\addcontentsline{toc}{chapter}{Annexe A --- Guide technique de déploiement et d'exécution}

\noindent Cette annexe décrit le guide technique de déploiement et d'exécution de SmartTrans-AI.

\section*{A.1 Backend}

\subsection*{Prérequis}

\begin{itemize}
  \item Java 25 ;
  \item Maven Wrapper fourni dans le dépôt ;
  \item PostgreSQL, via Docker Compose en environnement de développement ou via une instance locale.
\end{itemize}

\subsection*{Lancement}

Depuis le répertoire \texttt{backend} :

\begin{verbatim}
./mvnw spring-boot:run
\end{verbatim}

\subsection*{Vérification de l'état de l'application}

L'état de santé de l'API peut être vérifié à l'adresse :

\url{http://127.0.0.1:8080/actuator/health}

\subsection*{Documentation de l'API}

La documentation OpenAPI/Swagger est accessible à :

\url{http://127.0.0.1:8080/swagger-ui.html}

L'API est accessible sous le préfixe \texttt{/api/v1}.

\subsection*{Compte administrateur de démonstration}

Un compte administrateur est disponible en environnement de développement :

\texttt{admin@smarttrans.ai}

Le mot de passe est défini dans la configuration locale de l'application.

\section*{A.2 Frontend}

Depuis le répertoire \texttt{frontend}, installer les dépendances :

\begin{verbatim}
npm install
\end{verbatim}

Puis lancer l'application :

\begin{verbatim}
npm run dev
\end{verbatim}

L'interface est accessible à l'adresse :

\url{http://127.0.0.1:5173}

Le proxy Vite relaie les requêtes \texttt{/api} vers le backend exécuté sur le port \texttt{8080}.

\section*{A.3 Stack Docker}

Le fichier \texttt{docker-compose.prod.yml} permet d'assembler les principaux services de l'application :

\begin{itemize}
  \item PostgreSQL ;
  \item l'API backend ;
  \item l'interface frontend.
\end{itemize}

Les informations sensibles, notamment les mots de passe et les clés secrètes, doivent être configurées à l'aide d'un fichier d'environnement non versionné.
